\section{Neural Compute, Memory y Personalization}

Tras la demostración de acciones, la clase vuelve al nivel de la arquitectura para explicar qué lugar ocupa el modelo dentro del sistema y cómo el agent conserva información entre pasos y entre sesiones. Lo que más se presta a confusión aquí son context, persistent memory y user profile: los tres le proporcionan información al modelo, pero difieren en ciclo de vida, confiabilidad y obligaciones de privacidad.

\subsection{Model as Neural Compute Unit}

La slide Neural Compute Unit dibuja el Transformer como el núcleo de cómputo entre la entrada y la salida, y usa la CPU analogy para subrayar que un sistema más grande puede invocarlo. La analogía ayuda a entender la modularity, pero una CPU ejecuta instrucciones deterministas, mientras que un modelo de lenguaje genera distribuciones de probabilidad; por eso el agent también necesita reglas externas y un verifier que acoten las salidas inciertas.

\lecturefigure{slide-51-neural-compute-unit.jpg}{Analogía de la clase: el AI model como neural compute unit}{Video de la clase de Stanford 00:35:15}

\begin{knowledgebox}{Lectura de la figura: qué falta además de las flechas de entrada y salida}
En la figura, el model recibe instructions y context, y devuelve una response. Un agent real necesita además un scheduler que decida cuándo invocarlo, una memory layer que elija qué contexto usar, una tool layer que ejecute las acciones, una policy layer que verifique los permisos y un verifier que evalúe los resultados. El modelo es una compute unit, no un operating system completo.
\end{knowledgebox}

\teachervoice{El ponente recurre una y otra vez a la analogía CPU/chip: un modelo potente es solo una unidad de cómputo poderosa; lo que vuelve utilizable al sistema es la abstraction, la memory, la interface y el internet access que lo rodean. Esta idea atraviesa la figura del LLM OS que aparece más adelante.}

\subsection{Context y Long-term Memory}

La slide Long-term Memory divide la memory en memoria de trabajo y almacenamiento persistente. La context window se asemeja a un área de trabajo de corto plazo para la tarea actual y está limitada por el token budget; el persistent store guarda hechos entre sesiones, acciones pasadas o índices de documentos, y requiere gestionar la recuperación, la actualización, la eliminación y la provenance.

\lecturefigure{slide-52-long-term-memory.jpg}{Memoria de trabajo y memoria de largo plazo del agent}{Video de la clase de Stanford 00:35:33}

\begin{importantbox}{Flujo básico de la retrieval memory}
Para cada entrada de memoria $m_i$ se calcula su embedding $e_i$; si el vector de la consulta $q$ es $e_q$, las entradas pueden ordenarse según
\[
\operatorname{score}(m_i\mid q)
=\alpha\,\cos(e_q,e_i)
+\beta\,\operatorname{recency}(m_i)
+\gamma\,\operatorname{trust}(m_i)
\]
La similitud responde «¿es relevante?», recency se ocupa de la información desactualizada y trust registra la confiabilidad de la fuente. Los resultados recuperados todavía deben entrar en el context para que el modelo los use; la base de datos vectorial por sí misma no garantiza que los hechos sean correctos.
\end{importantbox}

\begin{warningbox}{El context no es memoria persistente, y los pesos del modelo tampoco son una base de datos de usuarios}
La ventana de contexto desaparece o se trunca al terminar la solicitud; la memoria persistente, en cambio, requiere un almacenamiento explícito. Escribir secretos del usuario en el prompt, en los registros o en una base de datos vectorial sin mecanismo de eliminación genera riesgos en el ciclo de vida de los datos. El conocimiento de entrenamiento contenido en los pesos del modelo tampoco puede cumplir la función de un registro personal que se pueda corregir con precisión.
\end{warningbox}

\teachervoice{En clase se hace la distinción de forma intuitiva: el context es la short-term memory que el modelo puede ver en ese momento, y solo una base de datos u otro store constituye la long-term memory. El ponente señala además que las tareas largas se topan rápidamente con el context limit, por lo que la memory management es un problema de sistema.}

\subsection{Personalization: preferencias explícitas, comportamiento implícito y retroalimentación}

La slide Personalization plantea que el user agent puede aprender preferences y usar las actions pasadas para mejorar su comportamiento posterior. Las preferencias explícitas provienen de lo que el usuario configura directamente, por ejemplo, su aerolínea preferida; las implícitas provienen de los registros de clics, selecciones y modificaciones. Estas últimas tienen más ruido y es más fácil que rebasen lo que el usuario espera. Esta sección parte de las fuentes de las señales para explicar cómo entran las preferencias en la toma de decisiones y por qué es indispensable permitir que el usuario las consulte, corrija y elimine.

\lecturefigure{slide-53-personalization.jpg}{Aprendizaje de preferencias personalizadas a partir de acciones históricas y retroalimentación}{Video de la clase de Stanford 00:37:39}

\begin{knowledgebox}{Términos clave: tres tipos de señales de personalización}
\begin{tabularx}{\textwidth}{L{0.18\textwidth}X X}
\toprule
Señal & Ejemplo & Principio de uso \\
\midrule
Explicit & «Solo vuelos directos», «No comprar automáticamente» & Se trata como restricción de alta prioridad; se puede consultar y editar \\
Implicit & Elegir varias veces cierta aerolínea o cierto modelo de calzado & Solo como soft preference; evitar inferencias excesivas a partir de una sola acción \\
Feedback & Aceptar, corregir, deshacer o presentar una queja & Registrar el contexto concreto de la tarea; distinguir un cambio de preferencia de un error del sistema \\
\bottomrule
\end{tabularx}
\end{knowledgebox}

Si la utilidad básica de la tarea para el candidato recomendado $x$ es $U_{task}(x)$, la puntuación de preferencia personalizada es $U_{pref}(x)$ y el riesgo es $C_{risk}(x)$, puede escribirse
\[
U(x)=U_{task}(x)+\lambda U_{pref}(x)-\mu C_{risk}(x).
\]
$\lambda$ no debe ser tan grande que anule las hard constraints, y $\mu$ debe aumentar con el dinero, la privacidad y la irreversibilidad en juego. El objetivo de la Personalization es reducir las explicaciones repetidas, no reescribir los objetivos en lugar del usuario.

\teachervoice{El ponente pone como ejemplo las preferencias al comprar zapatos y boletos de avión, y plantea aprender de manera continua a partir de human feedback. Estas notas añaden los límites necesarios: un clic del usuario puede ser solo una concesión momentánea y no debe fijarse para siempre como preferencia; el agent debe permitir consultar, corregir y olvidar.}

\subsection{Challenges: el aprendizaje continuo trae consigo problemas de privacidad y control}

La slide Challenges reúne collecting user preferences, active-learning preferences, learning from interactions, on-the-fly adaptation y privacy. Muestra que la personalization no consiste solo en mejorar la precisión de las recomendaciones, sino que es un problema de sistema integral que abarca la recolección de datos, la actualización en línea y el control por parte del usuario.

\lecturefigure{slide-54-agent-challenges.jpg}{Los retos centrales de la long-term personalization}{Video de la clase de Stanford 00:39:54}

\begin{warningbox}{Cuatro frentes de falla de un sistema de personalización}
\begin{enumerate}
  \item \textbf{Memoria errónea}: una sola operación equivocada se guarda a largo plazo;
  \item \textbf{Inferencia no autorizada}: se infieren atributos sensibles a partir del comportamiento o se reutilizan en otros contextos;
  \item \textbf{Preferencias obsoletas}: el usuario cambió, pero el sistema sigue actuando según el profile anterior;
  \item \textbf{Automatización irreversible}: el modelo convierte directamente una soft preference en una compra o en un envío de información hacia el exterior.
\end{enumerate}
Las medidas de mitigación incluyen provenance, políticas de caducidad, edición visible para el usuario, purpose limitation y confirmation para las acciones de alto riesgo.
\end{warningbox}

\teachervoice{Al hablar del futuro de la personalization, el ponente plantea por iniciativa propia las privacy concerns y considera las autonomous irreversible actions un punto de peligro. En clase no se trató «saber más del usuario» como un objetivo incondicional, sino que se discutió junto con los permisos y el control.}

\subsection{Resumen de la sección}

La analogía de la neural compute unit deja claro que el modelo es solo el núcleo de cómputo; el context aporta un área de trabajo de corto plazo, el persistent store aporta estado entre sesiones y la personalization convierte las preferencias explícitas, el comportamiento implícito y la retroalimentación en restricciones controlables. Cuanto más duradera es la memoria, más necesita mecanismos de provenance, actualización, eliminación, permisos y confirmación.
